```latex
We pass to the rotated azimuthal variable
\[
\psi=\phi-\frac{\pi}{4}.
\]
By abuse of notation, we write \(p_i\) both for the listed spherical coordinate pairs and for the corresponding points of \(\uS^2\). Let
\[
z=z(\theta,\psi)=(x,y,z_{\mathrm c})\in \uS^2,
\qquad
x=\sin\theta\cos\psi,\quad y=\sin\theta\sin\psi,\quad z_{\mathrm c}=\cos\theta.
\]
Whenever \(G_q(\theta,\psi)\) appears below, it means \(G_q(z(\theta,\psi))\). If
\[
(X,Y,Z)=(\sin\theta\cos\phi,\sin\theta\sin\phi,\cos\theta)
\]
are the usual Cartesian coordinates, then
\[
(X,Y,Z)=\left(\frac{x-y}{\sqrt2},\,\frac{x+y}{\sqrt2},\,z_{\mathrm c}\right),
\]
so \((x,y,z_{\mathrm c})\) is the Cartesian frame obtained by rotating the equatorial plane by \(-\pi/4\). In particular,
\[
\partial_\psi x=-y,\qquad \partial_\psi y=x,\qquad \partial_\psi z_{\mathrm c}=0,
\]
and, since \(\psi=\phi-\pi/4\), one also has \(\partial_\psi=\partial_\phi\).

\smallskip
\noindent\emph{(a)} Let
\[
\theta_0=\arccos\!\frac{1}{\sqrt3},
\qquad
\sin\theta_0=\sqrt{1-\frac13}=\sqrt{\frac23}.
\]
For \(\phi\in\{0,\pi/2,\pi,3\pi/2\}\), the rotated angle is
\[
\psi\in\left\{-\frac{\pi}{4},\,\frac{\pi}{4},\,\frac{3\pi}{4},\,\frac{5\pi}{4}\right\},
\]
hence
\[
(\cos\psi,\sin\psi)\in
\left\{
\frac{1}{\sqrt2}(1,-1),\,
\frac{1}{\sqrt2}(1,1),\,
\frac{1}{\sqrt2}(-1,1),\,
\frac{1}{\sqrt2}(-1,-1)
\right\}.
\]
Therefore, when \(\theta=\theta_0\), the corresponding rotated Cartesian coordinates are
\[
(x,y,z_{\mathrm c})\in
\left\{
\frac{1}{\sqrt3}(1,-1,1),\,
\frac{1}{\sqrt3}(1,1,1),\,
\frac{1}{\sqrt3}(-1,1,1),\,
\frac{1}{\sqrt3}(-1,-1,1)
\right\},
\]
and when \(\theta=\pi-\theta_0\), the same four pairs \((x,y)\) occur, while \(z_{\mathrm c}=-1/\sqrt3\). Writing
\[
q_{\varepsilon_1,\varepsilon_2,\varepsilon_3}:=
\frac{1}{\sqrt3}(\varepsilon_1,\varepsilon_2,\varepsilon_3),
\qquad
\varepsilon_j\in\{\pm1\},
\]
we obtain
\[
\begin{array}{llll}
p_1=q_{+,-,+}, & p_2=q_{+,+,+}, & p_3=q_{-,+,+}, & p_4=q_{-,-,+},\\[0.3em]
p_5=q_{+,-,-}, & p_6=q_{+,+,-}, & p_7=q_{-,+,-}, & p_8=q_{-,-,-}.
\end{array}
\]
These are exactly the eight points \(\frac{1}{\sqrt3}(\pm1,\pm1,\pm1)\), i.e. the vertices of the inscribed cube.

\smallskip
\noindent\emph{(b)} From the list in \emph{(a)},
\[
A=\{p_1,p_2,p_7,p_8\}
=\{q_{+,-,+},\,q_{+,+,+},\,q_{-,+,-},\,q_{-,-,-}\}.
\]
For each of these four sign triples one has \(\varepsilon_1\varepsilon_3=1\), and these are exactly the four sign triples with that property. Hence
\[
A=\{q_{\varepsilon_1,\varepsilon_2,\varepsilon_3}:\varepsilon_1\varepsilon_3=1\}.
\]
Likewise,
\[
B=\{p_3,p_4,p_5,p_6\}
=\{q_{\varepsilon_1,\varepsilon_2,\varepsilon_3}:\varepsilon_1\varepsilon_3=-1\}.
\]
Thus the \(+\) and \(-\) index sets of the \(T_{2g}\) generator are distinguished exactly by the sign of \(\varepsilon_1\varepsilon_3\).

\smallskip
\noindent\emph{(c)} Fix
\[
q=q_{\varepsilon_1,\varepsilon_2,\varepsilon_3}
=\frac{1}{\sqrt3}(\varepsilon_1,\varepsilon_2,\varepsilon_3),
\qquad
z=(x,y,z_{\mathrm c})\in \uS^2\setminus\{q\}.
\]
Since \(z\) and \(q\) are unit vectors,
\[
\cos(\dist(z,q))=z\cdot q
=\frac{\varepsilon_1 x+\varepsilon_2 y+\varepsilon_3 z_{\mathrm c}}{\sqrt3}.
\]
Therefore the half-angle identity gives
\[
\sin^2\!\frac{\dist(z,q)}{2}
=\frac{1-\cos(\dist(z,q))}{2}
=\frac{1-z\cdot q}{2}.
\]
Because \(0\le \dist(z,q)\le \pi\), one has \(\sin(\dist(z,q)/2)\ge 0\), so
\[
\log\!\sin\frac{\dist(z,q)}{2}
=\frac12\log(1-z\cdot q)-\frac12\log 2.
\]
Substituting this into
\[
G_q(z)=-\frac{1}{2\pi}\log\!\sin\frac{\dist(z,q)}{2}
\]
yields
\[
G_q(z)
=-\frac{1}{4\pi}\log(1-z\cdot q)+\frac{\log 2}{4\pi}.
\]
Thus
\[
G_q(z)=-\frac{1}{4\pi}\log(1-z\cdot q)+\mathrm{const},
\]
with \(\mathrm{const}=(\log 2)/(4\pi)\), independent of \(q\) and \(z\).

Now
\[
z\cdot q=\frac{\varepsilon_1 x+\varepsilon_2 y+\varepsilon_3 z_{\mathrm c}}{\sqrt3},
\]
hence
\begin{align*}
\partial_\psi G_q
&=
\partial_\psi\!\left[
-\frac{1}{4\pi}
\log\!\left(
1-\frac{\varepsilon_1 x+\varepsilon_2 y+\varepsilon_3 z_{\mathrm c}}{\sqrt3}
\right)
\right] \\
&=
\frac{1}{4\pi}
\frac{\partial_\psi\!\left((\varepsilon_1 x+\varepsilon_2 y+\varepsilon_3 z_{\mathrm c})/\sqrt3\right)}
{1-(\varepsilon_1 x+\varepsilon_2 y+\varepsilon_3 z_{\mathrm c})/\sqrt3} \\
&=
\frac{1}{4\pi\sqrt3}
\frac{\varepsilon_1\partial_\psi x+\varepsilon_2\partial_\psi y}
{1-(\varepsilon_1 x+\varepsilon_2 y+\varepsilon_3 z_{\mathrm c})/\sqrt3} \\
&=
\frac{1}{4\pi\sqrt3}
\frac{\varepsilon_1(-y)+\varepsilon_2 x}
{1-(\varepsilon_1 x+\varepsilon_2 y+\varepsilon_3 z_{\mathrm c})/\sqrt3} \\
&=
\frac{1}{4\pi\sqrt3}
\frac{\varepsilon_2 x-\varepsilon_1 y}
{1-(\varepsilon_1 x+\varepsilon_2 y+\varepsilon_3 z_{\mathrm c})/\sqrt3}.
\end{align*}
Summing over the four vertices in \(A\), and using \emph{(b)}, we get
\[
F_A(\theta,\psi)
:=
\partial_\psi\!\left(\sum_{q\in A} G_q(\theta,\psi)\right)
=
\frac{1}{4\pi\sqrt3}
\sum_{\substack{\varepsilon\in\{\pm1\}^3\\ \varepsilon_1\varepsilon_3=1}}
\frac{\varepsilon_2 x-\varepsilon_1 y}
{1-(\varepsilon_1 x+\varepsilon_2 y+\varepsilon_3 z_{\mathrm c})/\sqrt3}.
\]

\smallskip
\noindent\emph{(d)} By the same computation, now summing over \(B\), one has
\[
F_B(\theta,\psi)
=
\frac{1}{4\pi\sqrt3}
\sum_{\substack{\varepsilon\in\{\pm1\}^3\\ \varepsilon_1\varepsilon_3=-1}}
\frac{\varepsilon_2 x-\varepsilon_1 y}
{1-(\varepsilon_1 x+\varepsilon_2 y+\varepsilon_3 z_{\mathrm c})/\sqrt3}.
\]
On the other hand,
\[
x(\theta,\psi+\pi)=-x(\theta,\psi),\qquad
y(\theta,\psi+\pi)=-y(\theta,\psi),\qquad
z_{\mathrm c}(\theta,\psi+\pi)=z_{\mathrm c}(\theta,\psi).
\]
Hence
\begin{align*}
F_A(\theta,\psi+\pi)
&=
\frac{1}{4\pi\sqrt3}
\sum_{\substack{\varepsilon\in\{\pm1\}^3\\ \varepsilon_1\varepsilon_3=1}}
\frac{-\varepsilon_2 x+\varepsilon_1 y}
{1-(-\varepsilon_1 x-\varepsilon_2 y+\varepsilon_3 z_{\mathrm c})/\sqrt3}.
\end{align*}
Now make the substitution
\[
(\widetilde\varepsilon_1,\widetilde\varepsilon_2,\widetilde\varepsilon_3)
:=(-\varepsilon_1,-\varepsilon_2,\varepsilon_3).
\]
Then
\[
\widetilde\varepsilon_1\widetilde\varepsilon_3
=(-\varepsilon_1)\varepsilon_3
=-\varepsilon_1\varepsilon_3=-1,
\]
and the summand becomes
\[
\frac{\widetilde\varepsilon_2 x-\widetilde\varepsilon_1 y}
{1-(\widetilde\varepsilon_1 x+\widetilde\varepsilon_2 y+\widetilde\varepsilon_3 z_{\mathrm c})/\sqrt3}.
\]
Therefore
\[
F_A(\theta,\psi+\pi)
=
\frac{1}{4\pi\sqrt3}
\sum_{\substack{\widetilde\varepsilon\in\{\pm1\}^3\\ \widetilde\varepsilon_1\widetilde\varepsilon_3=-1}}
\frac{\widetilde\varepsilon_2 x-\widetilde\varepsilon_1 y}
{1-(\widetilde\varepsilon_1 x+\widetilde\varepsilon_2 y+\widetilde\varepsilon_3 z_{\mathrm c})/\sqrt3}
=
F_B(\theta,\psi).
\]

\smallskip
\noindent\emph{(e)} Define
\[
C(\theta,\psi):=F_A(\theta,\psi)+F_A(\theta,\psi+\pi),
\qquad
D(\theta,\psi):=F_A(\theta,\psi)-F_A(\theta,\psi+\pi).
\]
By \emph{(d)},
\[
C=F_A+F_B,
\qquad
D=F_A-F_B.
\]
Since \(A\) and \(B\) partition the eight sign triples according to \(\varepsilon_1\varepsilon_3=\pm1\), we obtain
\[
C(\theta,\psi)
=
\frac{1}{4\pi\sqrt3}
\sum_{\varepsilon\in\{\pm1\}^3}
\frac{\varepsilon_2 x-\varepsilon_1 y}
{1-(\varepsilon_1 x+\varepsilon_2 y+\varepsilon_3 z_{\mathrm c})/\sqrt3},
\]
and
\begin{align*}
D(\theta,\psi)
&=
\frac{1}{4\pi\sqrt3}
\left(
\sum_{\substack{\varepsilon\in\{\pm1\}^3\\ \varepsilon_1\varepsilon_3=1}}
-
\sum_{\substack{\varepsilon\in\{\pm1\}^3\\ \varepsilon_1\varepsilon_3=-1}}
\right)
\frac{\varepsilon_2 x-\varepsilon_1 y}
{1-(\varepsilon_1 x+\varepsilon_2 y+\varepsilon_3 z_{\mathrm c})/\sqrt3} \\
&=
\frac{1}{4\pi\sqrt3}
\sum_{\varepsilon\in\{\pm1\}^3}
\varepsilon_1\varepsilon_3\,
\frac{\varepsilon_2 x-\varepsilon_1 y}
{1-(\varepsilon_1 x+\varepsilon_2 y+\varepsilon_3 z_{\mathrm c})/\sqrt3}.
\end{align*}
This is exactly the claimed pair of formulas. By construction,
\[
C(\theta,\psi+\pi)=C(\theta,\psi),
\qquad
D(\theta,\psi+\pi)=-D(\theta,\psi).
\]

All identities in \emph{(c)}--\emph{(e)} are understood pointwise away from the relevant singular vertices.
```